% Options for packages loaded elsewhere
\PassOptionsToPackage{unicode}{hyperref}
\PassOptionsToPackage{hyphens}{url}
\PassOptionsToPackage{dvipsnames,svgnames,x11names}{xcolor}
\documentclass[
  11pt,
  table]{article}
\usepackage{xcolor}
\usepackage[paperwidth=8.5in,paperheight=8.5in,left=1.05in,right=1.05in,top=0.75in,bottom=0.8in]{geometry}
\usepackage{amsmath,amssymb}
\setcounter{secnumdepth}{-\maxdimen} % remove section numbering
\usepackage{iftex}
\ifPDFTeX
  \usepackage[T1]{fontenc}
  \usepackage[utf8]{inputenc}
  \usepackage{textcomp} % provide euro and other symbols
\else % if luatex or xetex
  \usepackage{unicode-math} % this also loads fontspec
  \defaultfontfeatures{Scale=MatchLowercase}
  \defaultfontfeatures[\rmfamily]{Ligatures=TeX,Scale=1}
\fi
\usepackage{lmodern}
\ifPDFTeX\else
  % xetex/luatex font selection
\fi
% Use upquote if available, for straight quotes in verbatim environments
\IfFileExists{upquote.sty}{\usepackage{upquote}}{}
\IfFileExists{microtype.sty}{% use microtype if available
  \usepackage[]{microtype}
  \UseMicrotypeSet[protrusion]{basicmath} % disable protrusion for tt fonts
}{}
\usepackage{setspace}
\makeatletter
\@ifundefined{KOMAClassName}{% if non-KOMA class
  \IfFileExists{parskip.sty}{%
    \usepackage{parskip}
  }{% else
    \setlength{\parindent}{0pt}
    \setlength{\parskip}{6pt plus 2pt minus 1pt}}
}{% if KOMA class
  \KOMAoptions{parskip=half}}
\makeatother
\usepackage{longtable,booktabs,array}
\usepackage{caption}
\captionsetup[table]{skip=6pt}
\newcounter{none} % for unnumbered tables
\usepackage{calc} % for calculating minipage widths
% Correct order of tables after \paragraph or \subparagraph
\usepackage{etoolbox}
\makeatletter
\patchcmd\longtable{\par}{\if@noskipsec\mbox{}\fi\par}{}{}
\makeatother
% Allow footnotes in longtable head/foot
\IfFileExists{footnotehyper.sty}{\usepackage{footnotehyper}}{\usepackage{footnote}}
\makesavenoteenv{longtable}
\usepackage{graphicx}
\makeatletter
\newsavebox\pandoc@box
\newcommand*\pandocbounded[1]{% scales image to fit in text height/width
  \sbox\pandoc@box{#1}%
  \Gscale@div\@tempa{\textheight}{\dimexpr\ht\pandoc@box+\dp\pandoc@box\relax}%
  \Gscale@div\@tempb{\linewidth}{\wd\pandoc@box}%
  \ifdim\@tempb\p@<\@tempa\p@\let\@tempa\@tempb\fi% select the smaller of both
  \ifdim\@tempa\p@<\p@\scalebox{\@tempa}{\usebox\pandoc@box}%
  \else\usebox{\pandoc@box}%
  \fi%
}
% Set default figure placement to htbp
\def\fps@figure{htbp}
\makeatother
\setlength{\emergencystretch}{3em} % prevent overfull lines
\providecommand{\tightlist}{%
  \setlength{\itemsep}{0pt}\setlength{\parskip}{0pt}}
\usepackage{fontspec}
\setmainfont{Helvetica Neue}[Ligatures=TeX,BoldFont={Helvetica Neue Bold},ItalicFont={Helvetica Neue Italic},BoldItalicFont={Helvetica Neue Bold Italic}]
\setsansfont{Helvetica Neue}[Ligatures=TeX,BoldFont={Helvetica Neue Bold},ItalicFont={Helvetica Neue Italic},BoldItalicFont={Helvetica Neue Bold Italic}]
\newfontfamily\fallbackfont{Arial Unicode MS}
\usepackage{newunicodechar}
\newunicodechar{₹}{Rs.\,}
\newunicodechar{→}{{\fallbackfont →}}
\newunicodechar{≈}{{\fallbackfont ≈}}
\newunicodechar{≥}{{\fallbackfont ≥}}
\newunicodechar{×}{{\fallbackfont ×}}
\usepackage{etoolbox}
\usepackage[most]{tcolorbox}
\usepackage{titlesec}
\usepackage{caption}
\usepackage{float}
\usepackage{adjustbox}
\usepackage{booktabs}
\usepackage{array}
\usepackage{fancyhdr}
\usepackage{tikz}
\definecolor{ink}{HTML}{22262B}
\definecolor{soft}{HTML}{6B7480}
\definecolor{paper}{HTML}{FFFFFF}
\definecolor{cPuzzle}{HTML}{E65100}
\definecolor{cContext}{HTML}{1565C0}
\definecolor{cPeople}{HTML}{6A1B9A}
\definecolor{cProcess}{HTML}{00838F}
\definecolor{cSystem}{HTML}{00695C}
\definecolor{cCause}{HTML}{AD1457}
\definecolor{cDesign}{HTML}{283593}
\definecolor{cTest}{HTML}{EF6C00}
\definecolor{cLook}{HTML}{455A64}
\definecolor{cWrong}{HTML}{C62828}
\definecolor{cOpen}{HTML}{E08600}
\definecolor{cOk}{HTML}{2E7D32}
\AtBeginDocument{\color{ink}}
\setlength{\parindent}{0pt}\setlength{\parskip}{6.5pt plus 1pt}
\renewcommand{\tightlist}{\setlength{\itemsep}{2pt}\setlength{\parskip}{0pt}}
\setcounter{secnumdepth}{0}
\captionsetup{font={small,it},labelformat=empty,skip=4pt,justification=raggedright,singlelinecheck=false,labelfont={color=soft}}
\floatplacement{figure}{!htbp}
\renewcommand{\topfraction}{.9}\renewcommand{\bottomfraction}{.9}\renewcommand{\textfraction}{.1}\renewcommand{\floatpagefraction}{.7}
\setcounter{topnumber}{2}\setcounter{bottomnumber}{2}\setcounter{totalnumber}{4}
\titleformat{\section}{\color{ink}\Large\bfseries}{}{0pt}{}
\titlespacing*{\section}{0pt}{16pt}{6pt}
\titleformat{\subsection}{\color{soft}\large\bfseries}{}{0pt}{}
\titlespacing*{\subsection}{0pt}{12pt}{4pt}
\pagestyle{fancy}\fancyhf{}\renewcommand{\headrulewidth}{0pt}
\fancyfoot[C]{\footnotesize\color{soft}\thepage}
\fancypagestyle{plain}{\fancyhf{}\renewcommand{\headrulewidth}{0pt}\fancyfoot[C]{\footnotesize\color{soft}\thepage}}
\AtBeginEnvironment{longtable}{\footnotesize\renewcommand{\arraystretch}{1.2}}
\setlength{\tabcolsep}{4pt}

% ---- devices
\tcbset{boxrule=0pt,arc=3pt,left=9pt,right=9pt,top=7pt,bottom=7pt,before skip=9pt,after skip=9pt,fonttitle=\bfseries\small}
\newtcolorbox{expcard}{colback=cLook!6,colframe=cLook,toprule=3pt,title={EXPERIMENT},coltitle=cLook,colbacktitle=cLook!6,attach boxed title to top left={yshift=-2pt},boxed title style={colback=cLook!6,boxrule=0pt},enhanced}
\newtcolorbox{turning}{colback=cOpen!12,colframe=cOpen,leftrule=5pt,toprule=0pt,rightrule=0pt,bottomrule=0pt,title={TURNING POINT},coltitle=cOpen!85!black,colbacktitle=cOpen!12,boxed title style={colback=cOpen!12},attach boxed title to top left={yshift=-2pt},enhanced}
\newtcolorbox{wrongbox}{colback=cWrong!6,colframe=cWrong,leftrule=4pt,toprule=0pt,rightrule=0pt,bottomrule=0pt,title={WE WERE WRONG ABOUT},coltitle=cWrong,colbacktitle=cWrong!6,boxed title style={colback=cWrong!6},attach boxed title to top left={yshift=-2pt},enhanced}
\newtcolorbox{tradeoff}{colback=cDesign!5,colframe=cDesign,leftrule=4pt,toprule=0pt,rightrule=0pt,bottomrule=0pt,title={THE TRADE-OFF},coltitle=cDesign,colbacktitle=cDesign!5,boxed title style={colback=cDesign!5},attach boxed title to top left={yshift=-2pt},enhanced}
\newtcolorbox{plainwords}{colback=cOk!6,colframe=cOk,toprule=2pt,leftrule=0pt,rightrule=0pt,bottomrule=0pt,title={IN PLAIN WORDS},coltitle=cOk,colbacktitle=cOk!6,boxed title style={colback=cOk!6},attach boxed title to top left={yshift=-2pt},enhanced}
\newtcolorbox{decisionbox}{colback=ink!5,colframe=ink,leftrule=4pt,toprule=0pt,rightrule=0pt,bottomrule=0pt,title={DECISION},coltitle=ink,colbacktitle=ink!5,boxed title style={colback=ink!5},attach boxed title to top left={yshift=-2pt},enhanced}
\newcommand{\pullquote}[1]{\begin{center}\begin{minipage}{0.88\linewidth}\raggedright\color{ink}\fontsize{17}{22}\selectfont\bfseries #1\end{minipage}\end{center}\vspace{4pt}}
\newcommand{\evtag}[2]{\textcolor{#1}{\scriptsize\textbf{\MakeUppercase{#2}}}}
\newcommand{\chapterpage}[5]{%
  \clearpage\noindent
  \begin{tikzpicture}[remember picture,overlay]
    \fill[#2] (current page.north west) rectangle ([yshift=-2.45in]current page.north east);
    \node[anchor=north west,text=white,inner sep=0pt] at ([xshift=0.8in,yshift=-0.3in]current page.north west) {\fontsize{84}{84}\selectfont\bfseries #1};
    \node[anchor=south east,text=white,inner sep=0pt] at ([xshift=-0.8in,yshift=-2.25in]current page.north east) {\fontsize{10}{12}\selectfont\bfseries #4};
  \end{tikzpicture}\par
  \vspace*{1.85in}
  {\noindent\color{#2}\bfseries\fontsize{26}{30}\selectfont\spaceskip=0.28em plus 0.1em minus 0.05em #3\par}
  \vspace{8pt}
  {\noindent\fontsize{13}{18}\selectfont\color{ink} #5\par}
  \vspace{10pt}}
\newcommand{\decisionstrip}[2]{\par\vspace{4pt}\noindent{\color{soft}\rule{\linewidth}{0.4pt}}\par\noindent{\scriptsize\color{soft}\textbf{DECISION #1} \quad #2}\par\vspace{6pt}}

% ---- cover and contents
\usepackage{tocloft}
\setlength{\cftbeforesecskip}{1.5pt}
\renewcommand{\cftsecfont}{\small}\renewcommand{\cftsecpagefont}{\small}
\renewcommand{\cfttoctitlefont}{\Large\bfseries\color{ink}}
\renewcommand{\maketitle}{%
  \thispagestyle{empty}
  \begin{tikzpicture}[remember picture,overlay]
    \fill[cLook!92!black] (current page.north west) rectangle (current page.south east);
    \foreach \i in {0,...,99} {
      \pgfmathtruncatemacro{\r}{int(\i/10)}\pgfmathtruncatemacro{\c}{mod(\i,10)}
      \ifnum\i<37 \fill[cPuzzle] ([xshift=1.0in+\c*0.2in,yshift=-0.9in-\r*0.2in]current page.north west) circle (0.065in);
      \else \draw[white!60!cLook,line width=0.6pt] ([xshift=1.0in+\c*0.2in,yshift=-0.9in-\r*0.2in]current page.north west) circle (0.065in);
      \fi}
    \node[anchor=north west,text=white,inner sep=0pt,text width=6.3in] at ([xshift=1.0in,yshift=-3.55in]current page.north west) {\bfseries\fontsize{40}{44}\selectfont\spaceskip=0.28em plus 0.1em minus 0.05em The Breakfast Paradox};
    \node[anchor=north west,text=white,inner sep=0pt,text width=5.6in] at ([xshift=1.0in,yshift=-5.05in]current page.north west) {\fontsize{15}{21}\selectfont\spaceskip=0.28em plus 0.1em Why a hostel books breakfast it does not eat, and what we would change};
    \node[anchor=south west,text=white,inner sep=0pt,align=left] at ([xshift=1.0in,yshift=0.8in]current page.south west) {\fontsize{11}{15}\selectfont\bfseries Team Invictus\\[2pt]\fontsize{10}{13}\selectfont\mdseries Systems Thinking and Design, IIIT Hyderabad};
    \node[anchor=south east,text=white!75!cLook,inner sep=0pt,align=right] at ([xshift=-1.0in,yshift=0.8in]current page.south east) {\fontsize{9}{12}\selectfont A book of experiments};
  \end{tikzpicture}
  \clearpage}
\usepackage{bookmark}
\IfFileExists{xurl.sty}{\usepackage{xurl}}{} % add URL line breaks if available
\urlstyle{same}
% fallback for those not using the hyperref driver hyperxmp:
\makeatletter
\@ifundefined{xmpquote}{\newcommand{\xmpquote}[1]{#1}}{}
\makeatother
\hypersetup{
  pdftitle={The Breakfast Paradox},
  colorlinks=true,
  linkcolor={cLook},
  filecolor={Maroon},
  citecolor={Blue},
  urlcolor={Blue},
  pdfcreator={LaTeX via pandoc}}

\title{The Breakfast Paradox}
\usepackage{etoolbox}
\makeatletter
\providecommand{\subtitle}[1]{% add subtitle to \maketitle
  \apptocmd{\@title}{\par {\large #1 \par}}{}{}
}
\makeatother
\subtitle{A book of experiments: why a hostel books breakfast it does not eat, and what we would change}
\author{}
\date{}

\begin{document}
\maketitle

\renewcommand*\contentsname{Contents}
{
\hypersetup{linkcolor=cLook}
\setcounter{tocdepth}{1}
\tableofcontents
}
\setstretch{1.15}
\clearpage\phantomsection\addcontentsline{toc}{section}{0  How to read this book}
\chapterpage{0}{cLook}{How to read this book}{BEFORE YOU START}{In about an hour you will know the problem, what we did about it, what we got wrong, what we recommend, and what nobody knows yet.}

This is the story of one student project at IIIT Hyderabad. We were asked to pick a puzzle on campus and study it as a system. We picked the hostel breakfast. Students are booked for it and billed for it, and most of them never eat it. The food is cooked anyway.

We wrote this book because the four reports we handed in were written one at a time, and each had to stand alone. They repeat each other. They also hide the part we find most useful, which is how our thinking changed. Here we tell it in the order it happened, with the reasons.

\textbf{The answer in one paragraph.} Most breakfasts that are booked are not eaten. Nobody chooses to book them: students are booked by default. Each kitchen cooks to its own guess of how many will turn up. The gap between booked and eaten is measured every day, in two places, and the measurement never reaches the people who set the rules. So we recommend building those missing routes first, opening a simple Sunday skip on a fixed date, and tuning the kitchens alongside. We say clearly where we are guessing.

\textbf{How each chapter works.} Every chapter asks one question, says what we did, shows what we saw, names the trade-off we faced, and records the decision and why we made it. Six small boxes recur.

\begin{expcard}

\textbf{Experiment.} Something we tried, what happened, and what surprised us.

\end{expcard}

\begin{turning}

\textbf{Turning point.} A finding that overturned something we believed.

\end{turning}

\begin{wrongbox}

A belief we held, the fact that replaced it, and why.

\end{wrongbox}

\begin{tradeoff}

Two honest options and what each would have cost us.

\end{tradeoff}

\begin{plainwords}

A hard word, explained once, in everyday language.

\end{plainwords}

\begin{decisionbox}

A numbered choice. All of them are listed at the back.

\end{decisionbox}

\textbf{How sure are we?} Facts come with a small tag. \evtag{cOk}{confirmed} means we read it off a record or saw it ourselves. \evtag{cContext}{single-sourced} means one person told us. \evtag{cOpen}{candidate} means we worked it out by putting sources together. \evtag{cLook}{assumed} means it is a judgement. In the last chapters, \evtag{cDesign}{projected} means a model output, \evtag{cCause}{simulated} means how we expect someone to react, and \evtag{cTest}{planned} means a test we have not run.

\begin{figure}
\centering
\includegraphics[width=1\linewidth,height=0.65\textheight,keepaspectratio,alt={The route of this book, from the puzzle to the recommendation.}]{Assets/f-journey.png}
\caption{The route of this book, from the puzzle to the recommendation.}
\end{figure}

\clearpage\phantomsection\addcontentsline{toc}{section}{1  The puzzle we were given}
\chapterpage{1}{cPuzzle}{The puzzle we were given}{PART 1: THE PUZZLE}{What were we actually trying to explain?}

\pullquote{

Why would a student book a breakfast, pay for it, and not go?

}

The problem we were handed was a small story about a student called Rohan. His alarm goes off four times at 7:45. He worked late, so twenty extra minutes of sleep feels better than breakfast. The hostel kitchen serves from 7:30 to 9:30, classes start at 8:30, and barely a third of the hostel shows up. Trays of food go in the bin. By 11:15 Rohan has no energy and cannot focus.

\begin{figure}
\centering
\includegraphics[width=1\linewidth,height=0.65\textheight,keepaspectratio,alt={Rohan\textquotesingle s morning, as the opening story told it. Every panel is a claim we had not yet checked.}]{Assets/f-rohan.png}
\caption{Rohan\textquotesingle s morning, as the opening story told it. Every panel is a claim we had not yet checked.}
\end{figure}

It is a good story. It has one clear villain, the clock. Late nights collide with early classes, students skip, food is wasted, energy crashes, and a hungry student buys a second meal at the canteen.

\textbf{What we did.} We wrote the story down as a chain of causes and asked of every link: do we have evidence for this, or only for the story? Only the first link, late nights leading to skipped mornings, had real testimony behind it. The energy crash and the second meal at the canteen had none. Our first interview was with ourselves, because one of us is a hostel student who skips breakfast. That gave us the first honest voice, and also a warning: one person is not a pattern.

\begin{figure}
\centering
\includegraphics[width=1\linewidth,height=0.65\textheight,keepaspectratio,alt={Our first guess as a chain. A question mark marks a link we had no evidence for.}]{Assets/f-firstguess.png}
\caption{Our first guess as a chain. A question mark marks a link we had no evidence for.}
\end{figure}

\begin{tradeoff}

\textbf{Believe the story and start fast, or doubt it and start slow.} The story was neat and gave us a plan on day one. But a neat story that turns out wrong shapes every later question. We chose to doubt it.

\end{tradeoff}

\begin{decisionbox}

\textbf{Decision 1.} We treated the opening story as a set of hypotheses, not facts, and wrote down which ones had evidence.

\end{decisionbox}

\textbf{What changed in our mind.} We stopped asking "why do students sleep through breakfast?" and started asking a wider question: what is the whole arrangement that produces a booking nobody acts on? That is a systems question, and it meant we needed to see the arrangement before we explained it.

\begin{plainwords}

\textbf{System.} A set of parts that affect each other, so that what happens in one place shows up in another. We use the word for the whole breakfast arrangement: the booking portal, the kitchens, the rules, the people.

\end{plainwords}

\clearpage\phantomsection\addcontentsline{toc}{section}{2  Drawing the edge}
\chapterpage{2}{cContext}{Drawing the edge}{PART 2: LOOKING FROM OUTSIDE}{What is inside the breakfast system, and what is not?}

\pullquote{

If the edge is in the wrong place, every answer after it is wrong.

}

\textbf{What we did.} We drew a boundary. We started from what a student sees: a portal, a QR code, a counter. Then we asked who and what that depends on. Three layers sat inside our boundary. The first is the booking and billing layer: the dining portal, the cancel button, the monthly bill and the QR scan at the counter that records who ate. The second is the kitchen layer: the catering firm that runs each of the four halls, from planning to cooking to top-up to leftovers. The third is the governance layer: the bodies that set the menu, the booking rules and the contracts.

\begin{figure}
\centering
\includegraphics[width=1\linewidth,height=0.65\textheight,keepaspectratio,alt={The breakfast system: three layers inside, edge actors on the line, class timetable and rival food outside.}]{Assets/f-edge.png}
\caption{The breakfast system: three layers inside, edge actors on the line, class timetable and rival food outside.}
\end{figure}

Some actors sat on the line itself. Suppliers, a garbage contractor, a site team that carries complaints, a weekly food-quality inspection, and people who are not students but eat the same food. Outside the line we put the class timetable, which shapes mornings but is run by someone else, and the other places a student can eat.

\textbf{What we saw.} The rules were stranger than the puzzle. A student is booked automatically for every meal, on campus or not, holidays included. The only way out is five cancellations per meal type per month, each at least four days ahead. After that the charge stands, whatever happens (\evtag{cOk}{confirmed}). The bill follows the booking, not the plate. A "Skip Meal" button exists, but it only tells the kitchen and refunds nothing. So a student who knows in advance has a better option than the official one: resell the booking on a student WhatsApp group.

\begin{figure}
\centering
\includegraphics[width=1\linewidth,height=0.65\textheight,keepaspectratio,alt={Booked by default: the pre-ticked box, the five exits, and the button that returns nothing.}]{Assets/f-default.png}
\caption{Booked by default: the pre-ticked box, the five exits, and the button that returns nothing.}
\end{figure}

Four halls with four kinds of food are held together only by this one portal. Each hall is run by its own caterer.

\begin{tradeoff}

\textbf{Study the whole mess, or only breakfast.} The whole mess would give us more to say. Breakfast only would let us go deep. We chose breakfast, and kept lunch and dinner as a comparison: same portal, same rules, same caterer, different result.

\end{tradeoff}

\begin{decisionbox}

\textbf{Decision 2.} Scope the study to breakfast. Keep lunch and dinner outside the boundary as a control, to test whether the cause is the meal or the system.

\end{decisionbox}

\begin{decisionbox}

\textbf{Decision 3.} Say how sure we are about every fact, with four labels, and keep the labels in every later document.

\end{decisionbox}

\textbf{What changed in our mind.} The first question grew a second half. We now also asked why each kitchen guesses how many to cook when the data to calibrate it is already recorded somewhere.

\clearpage\phantomsection\addcontentsline{toc}{section}{3  Who is in the room}
\chapterpage{3}{cPeople}{Who is in the room}{PART 2: LOOKING FROM OUTSIDE}{Who takes part in this system, and who did we forget?}

\textbf{What we did.} We made a long list of everyone who touches a breakfast, then looked at it through seventeen different lenses: students, money, technology, waste, governance, the academic calendar, sleep, rival food, and others. We asked of each name whether they had a real part, a real stake, or only a part in our imagination.

We started with thirty-two candidates. We ended with twenty-four. Three did not survive evidence at all: a campus health office, senior students passing on habits, and parents paying per meal. Nothing connected any of them to breakfast skipping.

\begin{figure}
\centering
\includegraphics[width=1\linewidth,height=0.65\textheight,keepaspectratio,alt={The cast: twenty-four stakeholders in rings, closest to the student\textquotesingle s morning at the centre.}]{Assets/f-cast.png}
\caption{The cast: twenty-four stakeholders in rings, closest to the student\textquotesingle s morning at the centre.}
\end{figure}

\begin{turning}

\textbf{Our first map had the cooks missing.} In the first version, "vendors" were one grey box. When we looked at who actually decides how much food appears, the caterers were not one actor. They were three firms with three ways of working, plus the managers, supervisors, kitchen teams and tasting staff inside them. We rebuilt the map with an operator cluster at its centre. We also added a group we had not imagined: people outside the student booking who eat from the same pots.

\end{turning}

\begin{figure}
\centering
\includegraphics[width=1\linewidth,height=0.65\textheight,keepaspectratio,alt={What the first map missed, and what the second one showed.}]{Assets/f-missing.png}
\caption{What the first map missed, and what the second one showed.}
\end{figure}

\textbf{What we saw.} The most important fact about the cast was an absence. The people who control the class timetable have no relationship of any kind with the people who run the mess. Nobody has raised the clash, and nobody has tried. We read that as inattention, not refusal. That matters because inattention is cheaper to fix.

\begin{tradeoff}

\textbf{Keep every candidate, or cut to those with evidence.} A long list is generous and hard to reason about. A short list is clear but may drop someone who matters later. We cut to twenty-four and kept the three dropped lenses on record, so the call is reviewable.

\end{tradeoff}

\begin{decisionbox}

\textbf{Decision 4.} Treat the operators as their own cluster, with the three firms kept separate, and cut the cast to twenty-four on evidence.

\end{decisionbox}

\clearpage\phantomsection\addcontentsline{toc}{section}{4  Who can change what}
\chapterpage{4}{cPeople}{Who can change what}{PART 2: LOOKING FROM OUTSIDE}{Who has the power to change this, and who actually uses it?}

\textbf{What we did.} We asked a harder question than "who is important." For each decision in the system (the booking rule, the cut-off, the menu, the first batch of food, the class start, the contract) we asked who proposes it, who approves it, and who feels the consequences. Then we marked two separate things on every actor: power, meaning the ability to decide, and leverage, meaning how much a small push from them would change the outcome.

\begin{plainwords}

\textbf{Leverage.} A small action that causes a large change in the whole system. A tiny turn of a ship\textquotesingle s rudder is leverage. A large push on the bow is power without leverage.

\end{plainwords}

\begin{figure}
\centering
\includegraphics[width=1\linewidth,height=0.65\textheight,keepaspectratio,alt={Power and interest, placed by judgement. Where an actor sits matters more than the exact spot.}]{Assets/f-power-grid.png}
\caption{Power and interest, placed by judgement. Where an actor sits matters more than the exact spot.}
\end{figure}

\textbf{What we saw.} Power was split into pieces that did not touch. A committee sets the menu. An office runs the portal and its rules. A warden holds the vendor contracts. The academic side holds the timetable. Each holds one piece, and nobody holds the whole outcome of "booked, cooked, eaten."

Students were the people who bear the cost and hold the least power. The caterers were the surprise: they hold the most knowledge and the least authority. Each caterer decides how much to cook, but none of them can change the booking rule, the cut-off, the bill or the menu. A person who sees the problem clearly and cannot touch its cause is a very particular kind of position.

\begin{figure}
\centering
\includegraphics[width=1\linewidth,height=0.65\textheight,keepaspectratio,alt={Who knows the gap, and who sets the rules. The bridge between them does not exist.}]{Assets/f-knows-decides.png}
\caption{Who knows the gap, and who sets the rules. The bridge between them does not exist.}
\end{figure}

\begin{plainwords}

\textbf{Inattention, not refusal.} When two parts of a system never speak, either someone tried and was blocked, or nobody tried. They are different problems with different cures.

\end{plainwords}

\begin{tradeoff}

\textbf{Score each actor once, or score each decision.} A single score per actor is simple. But it is wrong, because the same actor can hold strong authority over one decision and none over the next. It also makes people look powerful or powerless in general, which they are not.

\end{tradeoff}

\begin{decisionbox}

\textbf{Decision 5.} Judge power one decision at a time, and say whether it is formal authority or observed influence.

\end{decisionbox}

\textbf{What changed in our mind.} We stopped looking for a bad actor. Everyone we met had a sensible reason for what they did. The trouble was in how their sensible actions connected, and we could not see that without following the food itself.

\clearpage\phantomsection\addcontentsline{toc}{section}{5  Following one breakfast}
\chapterpage{5}{cProcess}{Following one breakfast}{PART 2: LOOKING FROM OUTSIDE}{What really happens between a booking and a bin?}

\textbf{What we did.} We traced one breakfast, step by step. For the student half we interviewed five students and used our own first account, six voices in all. They described how they decide: whether to wake, whether to check the app, whether to go early, whether to cancel, whether to resell. We were careful to say that six people are not a population, and that one transcript had a labelling gap, so we counted four to five distinct people.

\begin{figure}
\centering
\includegraphics[width=1\linewidth,height=0.65\textheight,keepaspectratio,alt={One student\textquotesingle s morning, with every point where a decision is made. Most decisions happen in the last minute.}]{Assets/f-morning-strip.png}
\caption{One student\textquotesingle s morning, with every point where a decision is made. Most decisions happen in the last minute.}
\end{figure}

\begin{expcard}

\textbf{We asked six students what they do about a booking they will not use. Result:} four had sold or bought a booking on the WhatsApp group, and none had used Skip Meal to get any value. The student who explained it said the button "does not return the money." \textbf{Surprise:} the unofficial market works better than the official tool, so the formal tool is simply ignored.

\end{expcard}

We also found a hypothesis we had been carrying that the data did not support. We had assumed friends pull each other to breakfast. Six respondents gave the same quiet answer: a friend might make a morning pleasanter, but nobody skips because a friend skips.

That was the student half. The kitchen half was dark. We did not know how a booking became a quantity of food, because nobody on our side of the counter could see it. So we went to look.

\begin{turning}

\textbf{We visited the kitchens, and asked to see their records.} One of the three kitchens let us photograph its books page by page. The caterers also talked us through their day. What we found rewrote the middle of our picture.

\end{turning}

There is no forecast. We had assumed the institution sends each kitchen a count four days ahead and the kitchen cooks to it. Instead each caterer reads the closed booking count and applies its own discount. Kadamba cooks a fixed share of the count. Yuktāhār plans from its own record books and a repeating two-week menu. Bakul cooks from a weekly judgement by its manager, and food travels in from a kitchen several kilometres away.

\begin{figure}
\centering
\includegraphics[width=1\linewidth,height=0.65\textheight,keepaspectratio,alt={Three kitchens, three ways of deciding how much to cook.}]{Assets/f-three-ways.png}
\caption{Three kitchens, three ways of deciding how much to cook.}
\end{figure}

\begin{figure}
\centering
\includegraphics[width=0.88\linewidth,height=0.65\textheight,keepaspectratio,alt={The books at Yuktāhār: one day\textquotesingle s production sheets, showing what was issued, cooked, left over, and when each item ran out.}]{Assets/plate-planning.jpg}
\caption{The books at Yuktāhār: one day\textquotesingle s production sheets, showing what was issued, cooked, left over, and when each item ran out.}
\end{figure}

\begin{wrongbox}

\textbf{A forecast locked four days ahead decides what the kitchen cooks.} The four-day close is on the booking, not on buying. Kitchens buy on their own schedules and set quantities the day before.

\textbf{Leftovers are composted.} At Kadamba they go to a garbage contractor. Nothing is composted there.

\textbf{Twenty to thirty staff eat what students leave.} About thirty extra portions are cooked in advance for thirty to fifty people from outside the student booking, such as night-shift security, and about half of Kadamba\textquotesingle s own staff also eat from the same food.

\end{wrongbox}

\begin{tradeoff}

\textbf{Describe "the operator", or trace each one.} One clean "operator" is easier to write about. But the three were different in exactly the way that mattered, and merging them would have hidden the finding. We traced each end to end.

\end{tradeoff}

\begin{decisionbox}

\textbf{Decision 6.} Treat each caterer as its own system, and trace each from booking to bin.

\end{decisionbox}

\begin{decisionbox}

\textbf{Decision 7.} When field evidence contradicts a finding already written, change the finding everywhere. We did this in three passes, as new evidence arrived in batches, and we did not blend old and new.

\end{decisionbox}

\textbf{What changed in our mind.} The kitchens were not passive. They were measuring the gap, correcting for it, and absorbing the cost, all privately. That made the next question obvious: how big is the gap, really?

\clearpage\phantomsection\addcontentsline{toc}{section}{6  The day the numbers arrived}
\chapterpage{6}{cProcess}{The day the numbers arrived}{PART 3: SEEING THE SYSTEM}{How big is the gap, and when does the crowd actually come?}

\pullquote{

We had been arguing from memory. Then the scan records arrived.

}

\textbf{What we did.} We asked the dining team for what only they hold: a scan for every breakfast booking at Kadamba for the whole of April 2026, 32,162 rows. A second set came from Yuktāhār, where we photographed its attendance register for twenty days in September, breakfast and lunch side by side. Until then, everything we knew about the crowd came from six people recalling their own mornings.

\begin{expcard}

\textbf{We counted booked against eaten, day by day. Result:} of 32,162 breakfast bookings at Kadamba in April, 11,941 were scanned, which is 37.1 percent, and 62.9 percent were not (\evtag{cOk}{confirmed}). \textbf{Surprise:} the booking count barely moved. It sat near 1,070 every day, Sunday included, while the number who ate rose and fell with the weekday.

\end{expcard}

\begin{figure}
\centering
\includegraphics[width=1\linewidth,height=0.65\textheight,keepaspectratio,alt={Out of a hundred breakfast bookings at Kadamba, about thirty-seven are eaten. At Yuktāhār breakfast, about twenty-eight.}]{Assets/f-hundred.png}
\caption{Out of a hundred breakfast bookings at Kadamba, about thirty-seven are eaten. At Yuktāhār breakfast, about twenty-eight.}
\end{figure}

\begin{turning}

\textbf{The gap was twice what we had written.} Our first submission said about thirty percent of bookings go unused. Two students had even said only thirty percent turn up, and we had read that as an unsupported impression. The records show the reverse. About sixty-three percent of Kadamba\textquotesingle s breakfasts go unused, and about seventy-two percent at Yuktāhār. Thirty percent was a lunch-sized number.

\end{turning}

\textbf{What we saw.} The weekday matters. Turnout at Kadamba ran from 42.1 percent on Tuesdays down to 27.3 percent on Sundays, with the three lowest days of the month all Sundays. A booking made once and left standing cannot follow a decision made each morning.

\begin{figure}
\centering
\includegraphics[width=0.82\linewidth,height=0.65\textheight,keepaspectratio,alt={Share of bookings eaten, by weekday, April. The booking count behind each bar is about the same.}]{Assets/d19-weekday-turnout.png}
\caption{Share of bookings eaten, by weekday, April. The booking count behind each bar is about the same.}
\end{figure}

The crowd comes late. It opens with a small burst at 7:30, runs flat until about nine, then climbs steeply. In the last fifteen minutes before the 9:30 close, 26.5 percent of all scans happen, and the rate is still rising at 9:29. Another 9.2 percent arrive after the close, through a back door the staff do not turn away. The latest scan was at 10:07.

\begin{figure}
\centering
\includegraphics[width=1\linewidth,height=0.65\textheight,keepaspectratio,alt={When Kadamba\textquotesingle s breakfast diners arrive. The last quarter hour is the busiest of the whole morning.}]{Assets/diagram8-crowd-curve.png}
\caption{When Kadamba\textquotesingle s breakfast diners arrive. The last quarter hour is the busiest of the whole morning.}
\end{figure}

At Yuktāhār we could do something we could not do anywhere else: compare breakfast and lunch from the same book, on the same days, with the same cooks. Lunch ran between 56 and 89 percent. Breakfast ran between 17.7 and 40.2 percent. Same portal, same rules, same caterer, very different result.

\begin{figure}
\centering
\includegraphics[width=1\linewidth,height=0.65\textheight,keepaspectratio,alt={Yuktāhār, September: breakfast against lunch, day by day. The dashed line is what the site head believed breakfast was.}]{Assets/d20-yuktahaar-breakfast-vs-lunch.png}
\caption{Yuktāhār, September: breakfast against lunch, day by day. The dashed line is what the site head believed breakfast was.}
\end{figure}

That dashed line is a quiet finding. The head of that kitchen keeps the best records on campus, and still believes breakfast turnout is about fifty percent. His own register says about twenty-eight. Even people who hold the numbers misjudge the level.

We also checked a worry. Six students were moving in different directions: one eating more, one eating less, one not at all. Did the aggregate move too? It did not. Each week of April sat between 35 and 39 percent. Individuals vary and the total holds still.

\begin{tradeoff}

\textbf{Wait for more records, or write with what we have.} New data was arriving in batches, and every batch meant revising. Waiting would have been safer and never ended. We froze the evidence at a stated date, wrote it up, and listed what was missing: Bakul has no records at all, and no month other than April or September.

\end{tradeoff}

\begin{decisionbox}

\textbf{Decision 8.} Use lunch at the same kitchen as the control for breakfast. Where recall and records disagree, trust the records.

\end{decisionbox}

\begin{decisionbox}

\textbf{Decision 9.} Freeze the evidence, state the date, and list the gaps instead of waiting.

\end{decisionbox}

\textbf{What changed in our mind.} We had been treating the gap as a puzzle about students. It is also a puzzle about kitchens, because if only a third show up, what do the kitchens cook for? The next chapter puts everything on one page.

\clearpage\phantomsection\addcontentsline{toc}{section}{7  One page for the whole system}
\chapterpage{7}{cSystem}{One page for the whole system}{PART 3: SEEING THE SYSTEM}{Can we draw all of it at once, and what stands out when we do?}

\textbf{What we did.} We put actors, flows, resources and constraints on one map, and asked of every part: what does it receive, what does it pass on, and where does it stop? We included the things beyond the kitchen door that shape it: suppliers, waste, the class timetable, and the informal ways people get around the rules.

\begin{figure}
\centering
\includegraphics[width=1\linewidth,height=0.65\textheight,keepaspectratio,alt={The system as a river: bookings flow into three kitchens, through one crowded morning, and out through five exits.}]{Assets/f-river.png}
\caption{The system as a river: bookings flow into three kitchens, through one crowded morning, and out through five exits.}
\end{figure}

\textbf{What we saw, kitchen by kitchen.} This is the one place in the book where we describe each kitchen fully.

Kadamba, run by a firm that holds the contract by tender, cooks a first batch for about 70 percent of bookings, planned separately for vegetarian and non-vegetarian and portioned by a gram standard for each dish. A supervisor watches the counters and calls for more, ready in 10 to 15 minutes (\evtag{cContext}{single-sourced}). Nothing in the kitchen carries one day\textquotesingle s result into the next day\textquotesingle s batch.

Bakul cooks 70 to 80 percent of bookings depending on the item. Its manager resets the share each week from his reading of past attendance, and he will not go below about 70 percent even on a very quiet weekend, because his rule is that there should be no shortage "even if some wastage occurs" (\evtag{cContext}{single-sourced}). The food is cooked at a kitchen about 8 km away, some of it from about 3 AM, and driven in insulated boxes. Temperature survives the journey and texture does not: hot parathas soften in their own steam. A top-up is a phone call, cooking, and a drive of 20 to 40 minutes. Nobody knows the total time.

Yuktāhār uses no fixed ratio. Twice a day the whole team meets and plans from the record of the same menu day two weeks earlier, a weekday attendance register, the month\textquotesingle s pattern and any events somebody mentions. During service it treats the first half hour as a sample of the rest. It weighs and records waste, and reuses leftovers: idli becomes upma, batter goes back in the fridge, milk becomes curd.

\begin{figure}
\centering
\includegraphics[width=1\linewidth,height=0.65\textheight,keepaspectratio,alt={The operators\textquotesingle{} mornings, side by side, with Bakul\textquotesingle s longer morning at the top.}]{Assets/d21-two-kitchen-daily-timeline.png}
\caption{The operators\textquotesingle{} mornings, side by side, with Bakul\textquotesingle s longer morning at the top.}
\end{figure}

Two things hold for all three kitchens. A cooked batch can be topped up but never shrunk, and the first batch is fixed before any diner arrives. At Kadamba, every meal passes a safety sequence first: cooking temperature, a sample kept for 72 hours, a wrapped display plate at the entrance, serving temperature and a tasting by the facility team. That sequence checks safety and taste. It never checks how much was cooked.

\begin{figure}
\centering
\includegraphics[width=1\linewidth,height=0.65\textheight,keepaspectratio,alt={Kadamba\textquotesingle s kitchen, seen on the visit: registers on the table, the cold room, and staff sorting vegetables.}]{Assets/plate-kitchen.jpg}
\caption{Kadamba\textquotesingle s kitchen, seen on the visit: registers on the table, the cold room, and staff sorting vegetables.}
\end{figure}

Where surplus goes is the other half of the map. Food is reused at Yuktāhār. At Kadamba it feeds the outside diners and half the staff, and some is held and reheated into later service, which we judge lowers its quality (\evtag{cOpen}{candidate}). Latecomers are served at the back door. What is left goes in the bin.

\begin{plainwords}

\textbf{Absorber.} Something that soaks up a consequence before anyone has to notice it. Resale absorbs a student\textquotesingle s loss. Outside diners absorb surplus food. Absorbers are usually kind, which is why they hide problems.

\end{plainwords}

One more thing the map showed. The system has already changed its own rules once, quickly. During a gas shortage, booking was made optional for a while and the cancellation limit was raised from five to ten (\evtag{cContext}{single-sourced}). Billing followed who came, and the no-show problem eased, on one senior account. Then the rules went back. The system knows how to do this. It does not know why it should keep doing it.

\begin{tradeoff}

\textbf{Draw everything, or draw something readable.} A map with every box is true and unreadable. We drew one clean picture around the booking count and the three kitchens, and kept the detail in the text.

\end{tradeoff}

\begin{decisionbox}

\textbf{Decision 10.} Build the map around the booking count, with three kitchens kept separate and every surplus exit shown, and keep academic timing on the map even though it has no channel to the mess.

\end{decisionbox}

\textbf{What changed in our mind.} Every piece of the map was reasonable. The kitchens adjust, the staff improvise, the students find workarounds. What the picture shows is a lot of local correction and no correction at the centre. That raised the obvious question: where does feedback actually work?

\clearpage\phantomsection\addcontentsline{toc}{section}{8  Where the loops close}
\chapterpage{8}{cSystem}{Where the loops close}{PART 3: SEEING THE SYSTEM}{Which of the feedback loops we suspected are real?}

\textbf{What we did.} We listed every feedback mechanism we had suspected, eighteen in all, and numbered them L1 to L17 plus one extra. For each, we set a strict test: a causal loop exists only if its last step really causes its first, and we have evidence for every link between them.

\begin{plainwords}

\textbf{Loop.} A chain of causes that bends back on itself. A \emph{balancing} loop pushes back against change, like a thermostat. A \emph{reinforcing} loop pushes the change further, like a rumour. A chain that never returns to its start is only a chain.

\end{plainwords}

\begin{figure}
\centering
\includegraphics[width=1\linewidth,height=0.65\textheight,keepaspectratio,alt={The eighteen mechanisms we tested, sorted by how much evidence each had.}]{Assets/f-loop-grid.png}
\caption{The eighteen mechanisms we tested, sorted by how much evidence each had.}
\end{figure}

\begin{expcard}

\textbf{We tried to close every loop. Result:} three closed with every link evidenced. L4 is how menu complaints reach the committee through the student council. L11 is Yuktāhār\textquotesingle s record-keeping, where today\textquotesingle s waste shapes the same menu day two weeks on. L12 is the counter top-up, which can only add food. Five more are candidates with one assumed link, and six are open chains. \textbf{Surprise:} every loop that really works sits inside a kitchen. None of them touches the booking.

\end{expcard}

\begin{figure}
\centering
\includegraphics[width=1\linewidth,height=0.65\textheight,keepaspectratio,alt={How Yuktāhār learns: what was left over shapes what is cooked, two weeks later.}]{Assets/f-l11.png}
\caption{How Yuktāhār learns: what was left over shapes what is cooked, two weeks later.}
\end{figure}

The most important chain did not close. We call it L13, the core of the paradox. Booking is automatic, so the count stays flat. The kitchen cooks a fixed share of a flat count. That share is too large on exactly the days and the vegetarian line where turnout is lowest. The surplus is held and reheated, which may lower quality. But nothing we found shows poor quality sending students away, so we could not close it.

\begin{figure}
\centering
\includegraphics[width=1\linewidth,height=0.65\textheight,keepaspectratio,alt={The core chain, link by link. The last arrow is dashed because we have no evidence for it.}]{Assets/f-chain.png}
\caption{The core chain, link by link. The last arrow is dashed because we have no evidence for it.}
\end{figure}

Resale is the clearest candidate loop. It is a shifting-the-burden pattern: relief for the student makes the problem tolerable, so nobody fixes the cause. Every link was evidenced except the last one, that relief reduces the pressure to fix billing.

\begin{wrongbox}

\textbf{Six of our loops were closed loops, named R1 and B1 to B4.} Three close (L4, L11, L12). The old names are retired. The shock rule from chapter 7, where booking went optional for a while, is not a loop either. It is a calendar switch that someone flips.

\end{wrongbox}

\begin{tradeoff}

\textbf{Call more things loops, or call fewer.} More loops makes the analysis look richer. But a loop with one assumed link is a hypothesis, and a report full of them cannot tell the reader what to trust. We kept the strict rule.

\end{tradeoff}

\begin{decisionbox}

\textbf{Decision 11.} Call something a loop only when every link is evidenced. Keep open chains as chains, with a dashed closing link.

\end{decisionbox}

\textbf{What changed in our mind.} The system is not short of feedback. It is short of it in the one place that matters. So the question became why the missing loops stay missing.

\clearpage\phantomsection\addcontentsline{toc}{section}{9  Why it keeps happening}
\chapterpage{9}{cCause}{Why it keeps happening}{PART 3: SEEING THE SYSTEM}{If everyone is sensible, why does the gap persist?}

\textbf{What we did.} We took seven concrete events and traced each downward, as far as it would go. This is the iceberg: we start with what you can see (an \emph{event}), look for the \emph{pattern} behind it, then the \emph{structure} that produces the pattern, and finally the belief, or \emph{mental model}, that holds the structure up.

\begin{figure}
\centering
\includegraphics[width=1\linewidth,height=0.65\textheight,keepaspectratio,alt={The paradox at four depths. The visible event sits on top. The belief that sustains it is at the bottom.}]{Assets/f-iceberg.png}
\caption{The paradox at four depths. The visible event sits on top. The belief that sustains it is at the bottom.}
\end{figure}

The seven chains covered: registering and not showing up; food running out before close while the bin fills; complaints that get no answer until something forces one; a fix that worked and was switched off; a menu change that helped and cast a shadow, because a known menu lets students decide in advance not to come; a real part of the gap that is just ordinary variation in life; and a kitchen that cooks to a guess and absorbs the difference.

\textbf{What we found.} Three structures recur in every chain. First, booking is an automatic standing default with a narrow, early exit, so it cannot follow a decision made in the morning. Second, billing fires at the booking and not at the plate, so an unused booking costs the system nothing. Third, the gap is measured twice, by the dining team\textquotesingle s scans and by the kitchens\textquotesingle{} own books, and neither measurement is routed to anyone who sets the rules. Two conditions keep all three in place. Attention goes to loud problems: a run-out has a complaint route that reaches a managing director, while a surplus is silent. And an absorption layer takes each consequence before it becomes a signal.

\begin{figure}
\centering
\includegraphics[width=1\linewidth,height=0.65\textheight,keepaspectratio,alt={Three structures and two conditions hold the gap in place.}]{Assets/f-pillars.png}
\caption{Three structures and two conditions hold the gap in place.}
\end{figure}

Underneath the kitchen side sits one mechanism, the conversion of bookings into a first batch, and it comes in three flavours. The caterers also hold three different attitudes. Kadamba cooks to the booking "because it is a responsibility and they are getting paid for it." Yuktāhār treats waste as a cost it wants to cut, and made its weighing blame-free before it worked. Bakul treats a shortage as the failure to avoid, and accepts waste as its price (\evtag{cContext}{single-sourced} for all three). The gap is about equally large at Kadamba and at Yuktāhār, so the cause sits in the booking. The kitchens\textquotesingle{} culture decides only how much of the gap turns into surplus (\evtag{cOpen}{candidate}).

Three patterns from systems thinking fit what we saw. \emph{Shifting the burden}: resale and kitchen skill each ease a symptom, so the cause is never touched. \emph{Fixes that fail}: a tightening of the cancellation rules, reportedly to cut office workload, may have strengthened the problem it aimed at (\evtag{cOpen}{candidate}, from secondhand reports). \emph{Accidental adversaries}: at Yuktāhār, weighing waste was first read as "they are catching us" and only worked once management made it blame-free.

\begin{wrongbox}

\textbf{Cuisine explains the low uptake at the north Indian halls.} We had read one portal snapshot as attendance. It showed bookings as a share of capacity, for one day, and did not match the kitchen\textquotesingle s own register. Inside Kadamba, vegetarian bookings are eaten at 32.4 percent and non-vegetarian at 48.5 percent, so the item and the line matter more than the cuisine.

\textbf{Every operator distrusts the booking.} Only Yuktāhār does. Kadamba treats the booking as an obligation, and Bakul adjusts by judgement.

\end{wrongbox}

One tension we could not resolve. Kadamba\textquotesingle s operator says about 5 to 10 percent of the food is wasted each day. But a first batch of 70 percent against turnout of 37 percent should leave far more. Either the outside diners, reheating and small top-ups absorb most of it, or the figure uses a different base. We do not know which. We left it open.

\begin{tradeoff}

\textbf{Find one root cause, or describe the structure.} A single root cause makes a clean headline and a clean fix. The evidence did not support one. We described structures, and said which one the others depend on.

\end{tradeoff}

\begin{decisionbox}

\textbf{Decision 12.} Read each chain down the iceberg separately, and put the cause in the booking structure, not in students or in operators.

\end{decisionbox}

\textbf{What changed in our mind.} We stopped asking who was at fault. The gap persists because it is comfortable for each person who touches it. The next step was to say that out loud as a new question.

\clearpage\phantomsection\addcontentsline{toc}{section}{10  Asking a better question}
\chapterpage{10}{cCause}{Asking a better question}{PART 3: SEEING THE SYSTEM}{What is the real problem, and what is the opportunity inside it?}

\pullquote{

The Breakfast Paradox is not a broken system. It is a stable one.

}

\textbf{What we did.} We reframed. Our first question was an operations question: why do students not eat what they book? We turned it into a systems statement. The system holds the gap between booking and eating in place because every actor\textquotesingle s behaviour inside it is reasonable, and every consequence is absorbed before it reaches anyone who sets the rules. Three structures hold it. Most of what follows is quietly absorbed, and most of the absorbers are useful, so none of the useful ones should be removed.

\begin{figure}
\centering
\includegraphics[width=1\linewidth,height=0.65\textheight,keepaspectratio,alt={Our question before and after the reframe.}]{Assets/f-fromto.png}
\caption{Our question before and after the reframe.}
\end{figure}

\begin{wrongbox}

\textbf{Five assumptions behind the first question did not survive.} Non-attendance is not always a problem: some skipping is ordinary. Students are not careless, and most are not choosing a booking at all. The fix is not unknown, because the institution has already run attendance-based billing, twice. The gap is not unmeasured: it is measured three ways and never reconciled. And it is not one problem.

\end{wrongbox}

We also tried other ways to frame it, and said why each was not enough. \emph{Student behaviour}: students skip for many reasons, and most did not choose to book. \emph{Operator waste}: operators already adjust, and cutting waste alone does not touch the booking. \emph{A clash with the 8:30 class}: late sleep is the reason students give most often, but the class start explains only part, and we cannot move it. \emph{A forecasting gap}: a better forecast would not reach the booking, and the kitchens\textquotesingle{} trouble is the batch they cannot shrink, not a missing prediction.

\textbf{The opportunity.} The system needs a way to notice its own cost. The gap is measured daily, and it persists because none of the measurements travels to anyone who can change the rules that produce it. So the work is to build missing routes, not to start measuring.

\begin{figure}
\centering
\includegraphics[width=1\linewidth,height=0.65\textheight,keepaspectratio,alt={Four places to act, laid along the booking-to-bin chain: the process, the rule, who decides, and how information travels.}]{Assets/f-ips.png}
\caption{Four places to act, laid along the booking-to-bin chain: the process, the rule, who decides, and how information travels.}
\end{figure}

The four places are these. On process, feed recorded turnout, by hall, weekday and line, into the decision that sets the first batch. On the rule, bring back a standing, breakfast-only version of the billing the institution has run twice. On decision rights, name one owner of the whole booked, cooked and eaten outcome. On information flow, add a trigger that fires on measured cost as well as on complaint volume.

We also wrote down what any intervention must respect. Do not remove an absorber to create a signal, and plan for the outside diners who eat the surplus. Act before the first batch is cooked, and never trade a safety step for speed. Cut waste without raising the chance of a shortage. Do not use turnout as the measure of success. Keep any new measurement blame-free, and build on the paper books kitchens already keep. Do not treat the class start as the lever, and do not aim the fix at students or at operators.

\begin{turning}

\textbf{One question stopped us.} What does a "plate" mean in each operator\textquotesingle s monthly payment: a booked plate or a served plate? Kadamba says it is paid for what is booked. Bakul\textquotesingle s manager says plates served. If it is booked plates, an operator earns on every no-show. If served plates, the operator pays for every uneaten portion. The answer decides who has a reason to close the gap, and nobody on our side has seen a contract. We marked it open and held back every idea that depended on it.

\end{turning}

\begin{tradeoff}

\textbf{Go for the deep fix, the booking, or the quick fix, the kitchen.} The booking is where the cause is, but it belongs to rule owners we cannot instruct. Kitchens are within reach but fixing them alone cannot move the student\textquotesingle s cost. We decided to plan both and keep them labelled as different things.

\end{tradeoff}

\begin{decisionbox}

\textbf{Decision 13.} Reframe as a stable equilibrium, held by three structures and two conditions.

\end{decisionbox}

\begin{decisionbox}

\textbf{Decision 14.} Never remove an absorber to create a signal, and name who eats each surplus.

\end{decisionbox}

\begin{decisionbox}

\textbf{Decision 15.} Hold every idea that depends on the payment basis until we know what a plate counts for.

\end{decisionbox}

\textbf{What changed in our mind.} We moved from describing to designing. The question was no longer "what is wrong?" but "where can a small change do the most, and what must it protect?"

\clearpage\phantomsection\addcontentsline{toc}{section}{11  Where we can push}
\chapterpage{11}{cDesign}{Where we can push}{PART 4: DECIDING WHAT TO CHANGE}{Where in this system would a small change do the most?}

\pullquote{

Changing a number is easy. Changing what the number is for is hard, and worth more.

}

\textbf{What we did.} We looked for places to push. Systems thinkers rank them by depth. Changing a \emph{parameter}, such as a ratio or a cut-off time, is shallow and easy. Changing a \emph{rule}, who gets to decide, how information flows, what the system is trying to do, or the belief it rests on, is deeper and harder. We worked with a twelve-place scale, counted from the deepest: place 1 is the power to step outside the whole way of thinking, place 2 is the mindset, place 3 the goal, place 5 the rules, place 6 the information flows, and place 12 the parameters.

\begin{plainwords}

\textbf{Leverage point.} A place in a system where a small, well-aimed change causes a large effect. They are rarely where you first look, and the deeper ones are usually the harder ones to move.

\end{plainwords}

We named twelve leverage points, LP1 to LP12, and laid them against the scale.

\begin{figure}
\centering
\includegraphics[width=1\linewidth,height=0.65\textheight,keepaspectratio,alt={Our twelve leverage points on the twelve-place depth scale. Held levers are ringed with a dashed line.}]{Assets/f-depth.png}
\caption{Our twelve leverage points on the twelve-place depth scale. Held levers are ringed with a dashed line.}
\end{figure}

LP1 is the enabler: build the missing routes, so the records that already exist reach the kitchens and the rule owners. LP2 is the deepest rule we can touch: change the breakfast default and its exit. LP3 to LP6 are the kitchen levers, which are an operator goal that treats surplus as cost, a forum where operators run trials together, a first batch calibrated to records, and a top-up that is aware of the crowd. LP7 is a trigger that fires on cost as well as on complaints. LP8 is a place for the operator\textquotesingle s voice in the menu. LP9 and LP10, a named owner of the outcome and what a booking costs, we held, because both wait on facts we do not have. LP11 and LP12, surfacing the beliefs and changing the goal, are the deepest.

We ruled out any lever that removes something protecting someone. We kept resale, reuse of leftovers, meals for the outside diners and staff, the walk-in buffer, and late service. We targeted only three things that protect nobody: the over-cooked first batch, repeated reheating, and the bin.

\textbf{Fourteen places to intervene.} Along the breakfast, from booking to bin, we marked fourteen points where a lever can bite, IP1 to IP14. They run from the default itself and the cancellation exit, through the day-before plan, the 5:30 first batch, the tasting, the first half hour and the 9:00 counter check, to the crowd at the close, the leftovers, the daily record, the bin, the records that go nowhere, and the menu cycle.

\textbf{Sunday is where three levers meet.} Kadamba\textquotesingle s Sunday turnout is 27.3 percent. A 70 percent vegetarian batch then cooks about 3.1 times what Sunday diners eat, against about 1.9 times on a Tuesday (\evtag{cOpen}{candidate}, our arithmetic). About 46 percent of Sunday scans arrive after 9:15. Three levers bite hardest there at once: the exit from the default, the calibrated batch, and the late batch. So Sunday on Kadamba\textquotesingle s vegetarian line became our first test slice.

\textbf{Why nothing changes by itself.} The present arrangement pays everyone something. A student gets no decision to make and a guaranteed seat. The caterers get the revenue they are used to, or protection from shortage. The institution gets a low workload and a predictable budget. We checked every lever against this: does it keep what each person values and change only what drives the over-booking? We also compared Kadamba with Yuktāhār, since they share the same rules yet one learns from its books and the other cooks a fixed ratio. A learning practice is therefore possible under today\textquotesingle s rules. But the two differ in five other ways, scale, lines, firm, month and measurement, so the comparison cannot prove why.

\textbf{Seven design principles.} These are the tests every idea must pass. We give them as short rules.

\begin{figure}
\centering
\includegraphics[width=1\linewidth,height=0.65\textheight,keepaspectratio,alt={The seven principles, as short rules.}]{Assets/f-principles.png}
\caption{The seven principles, as short rules.}
\end{figure}

\begin{tradeoff}

\textbf{Choose the deepest lever, or the one a student team can move.} The deepest levers are beliefs and goals, and a student team cannot instruct a committee. The ones we can move are routes, kitchen practice and a proposal. We ranked by depth, then judged feasibility honestly: could we prototype it, propose it to a forum that exists, or only write it down?

\end{tradeoff}

\begin{decisionbox}

\textbf{Decision 16.} Rank levers by depth, then judge each by what a student team can actually do. Say which levers are held and why.

\end{decisionbox}

\begin{decisionbox}

\textbf{Decision 17.} Make blame-free measurement a rule of the design. No recorded figure may be used against the person who recorded it.

\end{decisionbox}

\textbf{What changed in our mind.} We had expected one clever lever. The answer was a pairing: a deep fix that needs others, and a quick fix we can start now.

\clearpage\phantomsection\addcontentsline{toc}{section}{12  Trying ideas on paper}
\chapterpage{12}{cTest}{Trying ideas on paper}{PART 4: DECIDING WHAT TO CHANGE}{Which ideas survive when we test them against the whole system?}

\textbf{What we did.} We turned the levers into eight concepts and tested them. We could not pilot yet, so we tested on paper, with a model built on real data.

\begin{figure}
\centering
\includegraphics[width=1\linewidth,height=0.65\textheight,keepaspectratio,alt={The eight concepts, and the lever each one pulls. The Ratio Cut is the foil.}]{Assets/f-concepts.png}
\caption{The eight concepts, and the lever each one pulls. The Ratio Cut is the foil.}
\end{figure}

\textbf{A: The Breakfast Ledger.} One weekly sheet per hall that merges the scan records and the kitchen books, with the money charged for uneaten breakfasts. \textbf{B: The Kitchen Learning Kit.} A day-before plan card for Kadamba, with a named allowance for outside diners and staff, and a daily sheet. \textbf{C: The Crest Plan.} At 9:00 the supervisor starts a late batch that lands before the crowd. \textbf{D: The Sunday Exit.} A Sunday skip that actually removes the charge. \textbf{E: The Operator Menu Slot.} Operators bring what wastes and what runs out to the menu cycle, as a voice and not as authority. \textbf{F: Billed at Consumption.} Charge on the scan. \textbf{G: The Ratio Cut.} Cut Kadamba\textquotesingle s first batch to 50 percent by instruction. \textbf{H: The Shared Reading.} Joint sessions where each person reads their own number beside the record.

\begin{tradeoff}

\textbf{Test only ideas we believe in, or include one we expect to fail.} Testing only good ideas risks every test passing. We kept G in every test, on purpose, as a foil.

\end{tradeoff}

\begin{expcard}

\textbf{We cut Kadamba\textquotesingle s first batch to 50 percent, in the model. Result:} about 150 portions of surplus remain a day, and non-vegetarian run-outs begin (\evtag{cDesign}{projected}). \textbf{Surprise:} a blunt cut does not even fix the waste it targets. It moves the cost to the crowd at the close.

\end{expcard}

The model runs the proposed rules over twenty-three real April days of Kadamba scans. A calibrated first batch on its own cuts first-batch surplus from 362 to 60 portions a day, but runs 16.8 short at the crest. Adding a 9:00 late batch at 1.5 times cover, with a 10 percent buffer, gives 91 surplus, 2.6 short and 3.4 customers waiting (\evtag{cDesign}{projected}). Projecting at 8:00 instead of 9:00 is wrong by about 35 percent, so the projection waits until 9:00.

\begin{figure}
\centering
\includegraphics[width=0.88\linewidth,height=0.65\textheight,keepaspectratio,alt={Kitchen rules, backtested on April days. Each bar pair is surplus and shortage per day.}]{Assets/chart-a10-kitchen-scenarios.png}
\caption{Kitchen rules, backtested on April days. Each bar pair is surplus and shortage per day.}
\end{figure}

\begin{figure}
\centering
\includegraphics[width=0.88\linewidth,height=0.65\textheight,keepaspectratio,alt={The refined kitchen design, under shocks we assumed: a heavy day, an exam week, a favourite item.}]{Assets/chart-a10-stress-tests.png}
\caption{The refined kitchen design, under shocks we assumed: a heavy day, an exam week, a favourite item.}
\end{figure}

\textbf{Gates before scores.} We did not trust a single weighted score, so four gates make a concept ineligible whatever it scores. Do all absorbers survive? Is every safety step kept? Is it aimed at structure and not gated on an unknown that could reverse it? Is every kitchen lever paired with a route, default, goal or belief lever that lands on a fixed date? Concepts that pass are then scored 0 to 4 on nine criteria worth 400 points: booked against eaten, cooked against eaten, crest service, fundamental capability, fairness, food quality, desirability, feasibility and reversibility. The weights are our judgement, and we say so.

\begin{figure}
\centering
\includegraphics[width=1\linewidth,height=0.65\textheight,keepaspectratio,alt={Scores out of 100, with the gate each concept fails. Two bundles tie at the top.}]{Assets/f-scores.png}
\caption{Scores out of 100, with the gate each concept fails. Two bundles tie at the top.}
\end{figure}

The Ledger leads, because it builds the route every other lever needs, though on its own it moves no food and no money. The Kitchen Learning Kit alone has the biggest physical effect and the worst crest performance. Two bundles score the same, since a weighted score cannot see order. But one makes the Sunday skip wait until the kitchen has proved itself, and that fails the fourth gate. It is a small case of shifting the burden: if the kitchen fix comes first and works, nobody feels pressure to change the booking.

\textbf{How we expect people to react.} These are \evtag{cCause}{simulated}: reasoned from each person\textquotesingle s stated position, not actual feedback. Near-daily eaters would hardly notice. Sunday and frequent skippers would take a simple, free skip. A caterer paid on booked plates would oppose a skip, and one paid on served plates would be neutral. Kadamba\textquotesingle s supervisors fear a run-out and a reprimand, and would accept a late batch with a reserve and a named staff meal. Yuktāhār would welcome a paper ledger. The facilities department and the dining office are cautious about work they do not control. We wrote responses for twelve groups, and each one became a design change.

\begin{figure}
\centering
\includegraphics[width=1\linewidth,height=0.65\textheight,keepaspectratio,alt={Who we think would welcome, tolerate or resist each concept. These are expectations, not feedback.}]{Assets/f-responses.png}
\caption{Who we think would welcome, tolerate or resist each concept. These are expectations, not feedback.}
\end{figure}

\textbf{How it could go wrong.} Every second-order effect we found is a loop pushing back. A smaller batch could take food from the outside diners, so every plan card names an allowance and a missed meal stops the pilot. One loud shortage could bring the hedge back, so we judge over months, not days. A ledger can be read as an audit, which breeds round numbers and hidden leftovers, so it is blame-free and never used for a deduction. A grace period after 9:30 could soften the close, so it is withdrawn after four weeks if the late share stays above 9.2 percent.

\begin{decisionbox}

\textbf{Decision 18.} Keep a concept we expect to fail in every test, as a foil.

\end{decisionbox}

\begin{decisionbox}

\textbf{Decision 19.} Apply four gates before any score, and let the gates overrule the score.

\end{decisionbox}

\textbf{What changed in our mind.} The tests changed the design in twelve ways. The main change: the kitchen bundle must never run alone, and the exit must not wait for it.

\clearpage\phantomsection\addcontentsline{toc}{section}{13  What we recommend}
\chapterpage{13}{cTest}{What we recommend}{PART 4: DECIDING WHAT TO CHANGE}{What should actually happen, in what order, and who has to do it?}

\pullquote{

Route and read the gap together, open the Sunday exit on a fixed date, and tune the kitchen alongside.

}

\textbf{What we recommend.} The recommended design is three things that run together. The Ledger and the Shared Reading build the route. The Sunday skip changes the default on a fixed date, whatever the kitchen pilot shows. And the kitchen bundle, with its plan card, allowance, 9:00 late batch and daily sheet, runs beside them and is always labelled the quick fix.

\begin{figure}
\centering
\includegraphics[width=1\linewidth,height=0.65\textheight,keepaspectratio,alt={Twenty weeks in three stages. Decision points are on dates, not on results.}]{Assets/f-sequence.png}
\caption{Twenty weeks in three stages. Decision points are on dates, not on results.}
\end{figure}

The sequence has fixed dates. In weeks 0 to 4 we build the Ledger, run co-design tests, and hold a joint reading with the dining office and the kitchens. At week 6 the Sunday skip goes to the rule owner, with the plate-basis question attached. From week 5 to 8 the Kadamba Sunday vegetarian line runs the plan card and the late batch. If the skip is adopted it starts about week 10. At week 14 we decide whether an every-day skip follows, only if Sunday uptake is at least a quarter, no Sunday runs more than ten portions short, and the calibrated batch holds on both lines. Later steps, an opt-in on the quietest days and a change of billing, wait for facts we do not yet have. Handover begins at week 16, and we step back at week 20.

\textbf{The pilot.} Kadamba\textquotesingle s vegetarian line, Sunday breakfast, four Sundays, judged at week 16. The supervisor runs it, and we fill the ledger. It succeeds if Sunday vegetarian surplus falls from about 370 to under 100, no Sunday runs more than ten short, run-outs are no more frequent than on comparison days, every allowance meal is served, and grams and taste hold. It stops if shortage runs past one top-up two Sundays running, if a staff meal is missed, if portions fall below standard, if any figure is used against a person, or if consent is withdrawn.

\textbf{Who has to act.} Of thirteen actions, we can complete two ourselves, the Ledger and the plan cards, and the co-design tests. We can convene or support four that others decide. We can carry six forward as proposals with evidence: a quantity line at the tasting, a cost trigger, a stated goal, the Sunday exit, the billing basis and an owner of the outcome. And we must finish one ourselves: the handover. The team should not become the thing the system depends on. By week 20, two ledger pages and one forum should have run without us.

\textbf{What we expect.} All of this is \evtag{cDesign}{projected}. First-batch surplus at Kadamba falls from about 254 to 471 portions a day to about 69 to 120 over roughly four months. The weekly charge for uneaten breakfasts falls from about ₹2.47 lakh to ₹2.27 to ₹2.37 lakh with the Sunday skip, and to ₹1.23 to ₹1.85 lakh if the every-day skip follows. In April alone, ₹10.6 lakh was charged for uneaten Kadamba breakfasts. The kitchen change is large, and the change for students is small until the rule moves, which is what our diagnosis predicts.

\begin{figure}
\centering
\includegraphics[width=0.88\linewidth,height=0.65\textheight,keepaspectratio,alt={How Kadamba\textquotesingle s breakfast might respond over time, with ranges. Every line here is a projection.}]{Assets/chart-a10-bot-projection.png}
\caption{How Kadamba\textquotesingle s breakfast might respond over time, with ranges. Every line here is a projection.}
\end{figure}

\textbf{What stays out.} Bakul is not in the pilot. It has no records, and it is a different system: food cooked far away, a floor set by judgement, a payment basis reported as plates served. Nothing is sized for Bakul. The recommendation reaches it only once it has its own scans, a measured top-up time, an item-level record, and a confirmed payment basis.

\begin{tradeoff}

\textbf{Run the kitchen pilot first and let the results decide the exit, or fix the date.} Waiting for results feels careful. But a kitchen that gets better can make the booking look unimportant, and then the exit never happens. A fixed date keeps both moving.

\end{tradeoff}

\begin{decisionbox}

\textbf{Decision 20.} Put the Sunday skip on a fixed date, independent of how the kitchen pilot goes.

\end{decisionbox}

\begin{decisionbox}

\textbf{Decision 21.} Pilot on Sunday, Kadamba\textquotesingle s vegetarian line, where three levers meet and the effect is largest.

\end{decisionbox}

\begin{decisionbox}

\textbf{Decision 22.} Size nothing for Bakul until it has its own records.

\end{decisionbox}

\begin{decisionbox}

\textbf{Decision 23.} Plan our own exit. Hand the Ledger over to the dining office by role, and step back at week 20.

\end{decisionbox}

\textbf{What changed in our mind.} We started by wanting to explain why students skip breakfast. We ended with a design that asks very little of students and a good deal of the institution.

\clearpage\phantomsection\addcontentsline{toc}{section}{14  What we got wrong, and what nobody knows yet}
\chapterpage{14}{cLook}{What we got wrong, and what nobody knows yet}{PART 5: LOOKING BACK}{What would we tell someone who wants to continue this?}

\textbf{What we got wrong.} A quick index of the beliefs this project overturned, with the chapter that tells each story.

\begin{figure}
\centering
\includegraphics[width=0.7\linewidth,height=0.65\textheight,keepaspectratio,alt={Ten things we believed, and what the evidence said instead.}]{Assets/f-believed-found.png}
\caption{Ten things we believed, and what the evidence said instead.}
\end{figure}

\textbf{What we would do differently.} We would go to the kitchens first. For weeks we reasoned about the system from the student side and from a handbook, and the first honest look inside a kitchen changed more than any analysis. We would ask for the records sooner, because they existed all along. We would write down how sure we were from the very first document, which is the one habit that let us change our minds cleanly. And we would decide up front that when evidence overturns a finding, we change it everywhere, not in one place.

\textbf{What this taught us.} Measurement and routing are different things. This system measures its own problem every day, and it still cannot see it, because the measurement is not delivered. Reasonable people can produce an unreasonable outcome without anyone being at fault. And a fix that makes the visible symptom better can make the real problem harder to see.

\textbf{What nobody knows yet.}

\begin{figure}
\centering
\includegraphics[width=0.78\linewidth,height=0.65\textheight,keepaspectratio,alt={What is open, grouped by who could answer it.}]{Assets/f-open.png}
\caption{What is open, grouped by who could answer it.}
\end{figure}

A plate\textquotesingle s meaning in each operator\textquotesingle s payment. Who set the four-day close and the limit of five. Whether Skip Meal reaches Kadamba\textquotesingle s kitchen. Who heads the dining services body, since two different people are on record. What Kadamba\textquotesingle s waste figure is a percentage of. Whether a late batch can be cooked at 9:00 without slowing service. How many students would take a Sunday skip: we assumed a quarter to a half. Which booking rule applied to Yuktāhār in June. And everything about Bakul that a record would show. Every co-design test in this book is \evtag{cTest}{planned}, and none has run.

\textbf{The limits of our evidence.} We interviewed six students. We visited each kitchen once, and the caterers describe their own practice. Two sets of records cover two halls in two different months. The model uses only Kadamba\textquotesingle s scans. We have tried to put every one of those limits next to the claim it affects, not in a footnote.

\begin{figure}
\centering
\includegraphics[width=0.82\linewidth,height=0.65\textheight,keepaspectratio,alt={What the dining halls post on their own walls: the aim is to minimise food wastage through advance meal planning. This book is an attempt to show how the system could meet it.}]{Assets/plate-posters.jpg}
\caption{What the dining halls post on their own walls: the aim is to minimise food wastage through advance meal planning. This book is an attempt to show how the system could meet it.}
\end{figure}

\clearpage\phantomsection\addcontentsline{toc}{section}{Plain words}
\chapterpage{}{cLook}{Plain words}{BACK MATTER}{The terms in this book, in everyday language.}

\textbf{Absorber.} Something that soaks up a consequence before anyone has to notice it.

\textbf{Archetype.} A common pattern that shows up in many systems, such as shifting the burden: a quick fix eases the symptom, so the real cause is never addressed.

\textbf{Balancing loop.} A chain of causes that bends back and pushes against change, like a thermostat. \textbf{Reinforcing loop:} one that pushes the change further, like a rumour. \textbf{Chain:} causes that run in a line and never return to the start.

\textbf{Boundary.} The line that says what is inside the system we study and what is not.

\textbf{Delay.} The time between a cause and its effect. Long delays make a loop hard to see.

\textbf{Gate.} A pass-or-fail test a concept must clear before its score counts.

\textbf{Foil.} A weak idea kept in the test on purpose, to show how a good idea differs.

\textbf{Iceberg.} A way to read a problem at four depths: event, pattern, structure and belief.

\textbf{Leverage point.} A place where a small change makes a large difference.

\textbf{Stable equilibrium.} A state that resists change because the forces in it balance. A stable system can be a bad one.

\textbf{Projected, simulated, planned.} A model output, a reasoned expectation of someone\textquotesingle s reaction, and a test not yet run.

\section{What the codes mean}\label{what-the-codes-mean}

\textbf{The loops, L1 to L17.} L1 resale as a safety net. L2 Skip Meal. L3 students moving from one hall to another, and L3b lock-in on Kadamba. L4 menu rotation through the council and committee. L5 menu transparency. L6 the shock rule. L7 late nights leading to skipping. L8 vendor discipline by the market. L9 absorption of waste. L10 the measured but unrouted gap. L11 Yuktāhār\textquotesingle s record-keeping. L12 top-up at the counter. L13 the core chain from booking to surplus. L14 the closing rush and the just-in-case batch. L15 the soft close. L16 operator waste data reaching the menu. L17 kitchen skill shifting the burden.

\textbf{The leverage points, LP1 to LP12.} LP1 build the missing routes. LP2 change the default and its exit. LP3 operator goal. LP4 operator forum. LP5 calibrated first batch. LP6 crest-aware top-up and late batch. LP7 cost trigger. LP8 operator voice in the menu. LP9 owner of the outcome (held). LP10 what a booking costs (held). LP11 test the beliefs. LP12 change the goal.

\textbf{The principles, P1 to P7.} Route before you measure. Make measurement blame-free. Change the default, goal and belief, not the people. Act on the first batch the day before. Feed the absorbers by plan. Measure cooked against eaten, run-outs and quality, not turnout. Use plural measures, never one target.

\section{Decision log}\label{decision-log}

{\def\LTcaptype{none} % do not increment counter
\begin{longtable}[]{@{}
  >{\raggedright\arraybackslash}p{(\linewidth - 4\tabcolsep) * \real{0.0392}}
  >{\raggedright\arraybackslash}p{(\linewidth - 4\tabcolsep) * \real{0.8431}}
  >{\raggedright\arraybackslash}p{(\linewidth - 4\tabcolsep) * \real{0.1176}}@{}}
\toprule\noalign{}
\begin{minipage}[b]{\linewidth}\raggedright
\#
\end{minipage} & \begin{minipage}[b]{\linewidth}\raggedright
Decision
\end{minipage} & \begin{minipage}[b]{\linewidth}\raggedright
Chapter
\end{minipage} \\
\midrule\noalign{}
\endhead
\bottomrule\noalign{}
\endlastfoot
1 & Treat the opening story as hypotheses, not facts & 1 \\
2 & Study breakfast only, with lunch and dinner as a control & 2 \\
3 & Label how sure we are about every fact & 2 \\
4 & Treat the operators as a cluster of three, and cut the cast to twenty-four & 3 \\
5 & Judge power decision by decision & 4 \\
6 & Trace each caterer separately, from booking to bin & 5 \\
7 & When evidence overturns a finding, change it everywhere & 5 \\
8 & Use lunch as the control, and trust records over recall & 6 \\
9 & Freeze the evidence, state the date, list the gaps & 6 \\
10 & Draw the map around the booking count and three kitchens & 7 \\
11 & Call it a loop only if every link is evidenced & 8 \\
12 & Put the cause in the booking structure, not in people & 9 \\
13 & Reframe as a stable equilibrium & 10 \\
14 & Never remove an absorber to create a signal & 10 \\
15 & Hold every idea that depends on the payment basis & 10 \\
16 & Rank by depth, judge by feasibility, name what is held & 11 \\
17 & Blame-free measurement is a rule of the design & 11 \\
18 & Keep a concept we expect to fail, as a foil & 12 \\
19 & Gates before scores & 12 \\
20 & Sunday skip on a fixed date & 13 \\
21 & Pilot on Kadamba\textquotesingle s Sunday vegetarian line & 13 \\
22 & Size nothing for Bakul without records & 13 \\
23 & Plan our own exit at week 20 & 13 \\
\end{longtable}
}

\section{Where the original deliverables live}\label{where-the-original-deliverables-live}

{\def\LTcaptype{none} % do not increment counter
\begin{longtable}[]{@{}
  >{\raggedright\arraybackslash}p{(\linewidth - 4\tabcolsep) * \real{0.3438}}
  >{\raggedright\arraybackslash}p{(\linewidth - 4\tabcolsep) * \real{0.5469}}
  >{\raggedright\arraybackslash}p{(\linewidth - 4\tabcolsep) * \real{0.1094}}@{}}
\toprule\noalign{}
\begin{minipage}[b]{\linewidth}\raggedright
Deliverable
\end{minipage} & \begin{minipage}[b]{\linewidth}\raggedright
Its sections, in this book
\end{minipage} & \begin{minipage}[b]{\linewidth}\raggedright
Chapter
\end{minipage} \\
\midrule\noalign{}
\endhead
\bottomrule\noalign{}
\endlastfoot
System Context Brief & Boundary and context; purpose and intended outcomes; symptoms and underlying issues; actors, institutions, processes, resources, constraints; initial problem framing & 1, 2 \\
Stakeholder Map & Use case; primary, secondary and indirect stakeholders; interests and expectations; formal and informal roles; the maps; relationships and conflicts; findings & 3 \\
Power-Interest and Leverage Map & Decision-making authority; formal and informal power; incentives; who can and cannot influence; the map; implications & 4 \\
Process Trace & Interviews and observations; decisions and interactions; formal and informal rules; escalation; exceptions, workarounds, dependencies & 5 \\
System Timeline & Records and observations; timelines; significant changes; feedback loops and patterns; unintended consequences & 6, 8 \\
System Map & Actors, process, resources, constraints; three kitchens; waste; service timeline; conflicts; formal and informal structure; bottlenecks and tensions; what stands out & 7 \\
Causal Loop and Feedback Analysis & Dashboard; reinforcing and balancing loops; candidates; chains that do not close; absent mechanisms; unintended consequences; structural causes & 8 \\
Systemic Problem Analysis & Events, patterns, structures, mental models; structural causes; loops; archetypes; consequences; tensions; power; workarounds; synthesis & 9 \\
Reframed Problem Statement & Challenging assumptions; stakeholder perspectives; contradictions; the iceberg; why it holds steady; alternative framings; validation & 10 \\
Design Opportunity Statement & Where the leverage is; the opportunity; validation; four intervention points; what must be respected; how they interact; gating questions; what Phase 3 inherits & 10 \\
Leverage-Point Map and Design Principles & Leverage points; changes by type; the case for the status quo; intervention points; the map; principles P1 to P7 & 11 \\
Intervention Concepts, Scenario Model, Final Recommendation & Alternatives; scenarios and prototypes; simulated responses; tests against objectives; unintended consequences; refinement; recommendation, pilot, monitoring & 12, 13 \\
Open questions and limits & Outstanding data collection, in every deliverable & 14 \\
\end{longtable}
}

\section{What this book is built from}\label{what-this-book-is-built-from}

Six student interviews. Four wall posters. A talk with the head of facilities. Three caterer visits and 26 photographed pages of one kitchen\textquotesingle s books. Every Kadamba breakfast scan from April 2026. A summary of the Bakul talk.

\end{document}
